\documentclass[11pt]{article}
\usepackage{cornell-notes}
% Suite identity: Web Design Patterns / Design of Sites (document DOS-05).
% This study document publishes under the Cornell Notes tree; the callout
% environments and provenance marker below are document-local mirrors of the
% suite's reconstructed style (see
% src/programming/web/design-of-sites/design-of-sites-style.tex).
\usepackage{booktabs}
\usepackage{enumitem}
\definecolor{Ink}{HTML}{1E2A36}
\definecolor{Ocean}{HTML}{0F4C72}
\definecolor{Mint}{HTML}{0E7E6B}
\definecolor{Amber}{HTML}{B07A1E}
\definecolor{Faint}{HTML}{F2F7F9}
\newtcolorbox{groundedbox}[1][]{enhanced,breakable,colback=Faint,colframe=Mint,boxrule=0.5pt,arc=1.5mm,title={Grounded in the public record},fonttitle=\bfseries\small,#1}
\newtcolorbox{dosnotebox}[1][]{enhanced,breakable,colback=Faint,colframe=Amber,boxrule=0.5pt,arc=1.5mm,#1}
\newcommand{\Provenance}{\par\noindent{\footnotesize\sffamily\color{Amber}\bfseries Provenance:}~{\footnotesize\color{Ink}\itshape claims grounded in the public bibliographic record are boxed and cited; everything else is newly authored interpretation.}\par\smallskip}
\setCornellCollection{Web Design Patterns Cornell Notes}
\setCornellSourceTitle{The Design of Sites}
\setCornellSourceEdition{2nd edition}
\setCornellSourceAuthor{van Duyne, Landay and Hong}
\title{The Design of Sites: Cornell Study Notes}
\author{Web Design Patterns Engineering Suite (DOS-05)}
\date{}
\begin{document}
\maketitle
\section{How to study with these notes}

\Provenance
\begin{groundedbox}The five foundational chapter titles and groups A--M are verified in [B1]. The cue/notes/summary material is an \emph{original conceptual study aid inferred from those topics and contemporary sources}; it is \textbf{not a line-by-line summary of chapters whose full text was not retrieved}.\end{groundedbox}
\textbf{Cornell protocol.} Review the prompt in the left column without looking at notes; answer from memory; compare; complete the exercise; write a 1--2 sentence retrospective in your own words. Use the catalog [DOS-02] to deepen any pattern ID referenced here.

\section{Study Unit 01: Foundation 1 — Customer-Centered Design}
\begin{longtable}{@{}p{.27\linewidth}p{.66\linewidth}@{}}\toprule
\textbf{Cue / Recall} & \textbf{Notes / Relationships}\\\midrule\endhead
Whose goals matter, and how are they observed? & Frame the work as representative users completing meaningful tasks. Record user, business and technical constraints separately so the team can see when they conflict.\\
\addlinespace[4mm]
\textbf{Practice} & Write three behavioral hypotheses for a checkout, then design a study that can disprove them.\\\bottomrule\end{longtable}
\begin{dosnotebox}\textbf{Bottom summary.} User outcomes should control the design agenda rather than interface fashion.\\[2mm]
\textbf{Your own synthesis:} \rule{.72\linewidth}{.4pt}\\[4mm]
\rule{.96\linewidth}{.4pt}\end{dosnotebox}

\section{Study Unit 02: Foundation 2 — Pattern Languages}
\begin{longtable}{@{}p{.27\linewidth}p{.66\linewidth}@{}}\toprule
\textbf{Cue / Recall} & \textbf{Notes / Relationships}\\\midrule\endhead
When is a reusable pattern applicable? & A design pattern is a documented response to a recurring problem in a specific context. Record forces, alternatives, drawbacks and associated tests.\\
\addlinespace[4mm]
\textbf{Practice} & Compare K3 tabs with navigational links for two different journeys.\\\bottomrule\end{longtable}
\begin{dosnotebox}\textbf{Bottom summary.} Patterns are conditional decisions, not universal prescriptions.\\[2mm]
\textbf{Your own synthesis:} \rule{.72\linewidth}{.4pt}\\[4mm]
\rule{.96\linewidth}{.4pt}\end{dosnotebox}

\section{Study Unit 03: Foundation 3 — Knowing Customers}
\begin{longtable}{@{}p{.27\linewidth}p{.66\linewidth}@{}}\toprule
\textbf{Cue / Recall} & \textbf{Notes / Relationships}\\\midrule\endhead
What counts as credible research evidence? & Contextual inquiry, participant tasks, analytics and support cases each reveal different slices of reality. Bias and accessibility omissions should be logged.\\
\addlinespace[4mm]
\textbf{Practice} & Plan six research participants across device and ability variation.\\\bottomrule\end{longtable}
\begin{dosnotebox}\textbf{Bottom summary.} Triangulation is stronger than one research technique.\\[2mm]
\textbf{Your own synthesis:} \rule{.72\linewidth}{.4pt}\\[4mm]
\rule{.96\linewidth}{.4pt}\end{dosnotebox}

\section{Study Unit 04: Foundation 4 — Iterative Design}
\begin{longtable}{@{}p{.27\linewidth}p{.66\linewidth}@{}}\toprule
\textbf{Cue / Recall} & \textbf{Notes / Relationships}\\\midrule\endhead
How should uncertainty be reduced? & Start with hypotheses and inexpensive prototypes; test risky assumptions first. High fidelity does not compensate for invalid task organization.\\
\addlinespace[4mm]
\textbf{Practice} & Run two rounds of first-click testing and document design change.\\\bottomrule\end{longtable}
\begin{dosnotebox}\textbf{Bottom summary.} Iteration is an evidence loop, not a sequence of screens.\\[2mm]
\textbf{Your own synthesis:} \rule{.72\linewidth}{.4pt}\\[4mm]
\rule{.96\linewidth}{.4pt}\end{dosnotebox}

\section{Study Unit 05: Foundation 5 — Development Process}
\begin{longtable}{@{}p{.27\linewidth}p{.66\linewidth}@{}}\toprule
\textbf{Cue / Recall} & \textbf{Notes / Relationships}\\\midrule\endhead
How does research reach production? & Translate discoveries into acceptance criteria, architecture decisions, components and service contracts, then validate outcomes after deployment.\\
\addlinespace[4mm]
\textbf{Practice} & Create a definition of done combining usability, security and field metrics.\\\bottomrule\end{longtable}
\begin{dosnotebox}\textbf{Bottom summary.} A shipped screen is not proof that a task works.\\[2mm]
\textbf{Your own synthesis:} \rule{.72\linewidth}{.4pt}\\[4mm]
\rule{.96\linewidth}{.4pt}\end{dosnotebox}

\section{Study Unit 06: A — Site Genres}
\begin{longtable}{@{}p{.27\linewidth}p{.66\linewidth}@{}}\toprule
\textbf{Cue / Recall} & \textbf{Notes / Relationships}\\\midrule\endhead
Which kind of service is being built? & A service genre predicts its central user missions and trust expectations. Genre guides journeys but does not prescribe a backend topology.\\
\addlinespace[4mm]
\textbf{Practice} & Compare primary tasks in educational and government portals.\\\bottomrule\end{longtable}
\begin{dosnotebox}\textbf{Bottom summary.} Choose UX structure from mission and audience.\\[2mm]
\textbf{Your own synthesis:} \rule{.72\linewidth}{.4pt}\\[4mm]
\rule{.96\linewidth}{.4pt}\end{dosnotebox}

\section{Study Unit 07: B — Navigation Framework}
\begin{longtable}{@{}p{.27\linewidth}p{.66\linewidth}@{}}\toprule
\textbf{Cue / Recall} & \textbf{Notes / Relationships}\\\midrule\endhead
How can users find content from different mental models? & Task-based, categorical, alphabetical, chronological and popular access paths can coexist if taxonomy and route identity stay consistent.\\
\addlinespace[4mm]
\textbf{Practice} & Conduct a tree test for B3/B4 and validate site accessibility B9.\\\bottomrule\end{longtable}
\begin{dosnotebox}\textbf{Bottom summary.} Findability is a user and information-model property.\\[2mm]
\textbf{Your own synthesis:} \rule{.72\linewidth}{.4pt}\\[4mm]
\rule{.96\linewidth}{.4pt}\end{dosnotebox}

\section{Study Unit 08: C — Powerful Homepage}
\begin{longtable}{@{}p{.27\linewidth}p{.66\linewidth}@{}}\toprule
\textbf{Cue / Recall} & \textbf{Notes / Relationships}\\\midrule\endhead
What must the initial view accomplish? & The homepage can communicate purpose and route people to tasks. Avoid excessive content modules that delay comprehension and rendering.\\
\addlinespace[4mm]
\textbf{Practice} & Test whether users identify the core offer in five seconds.\\\bottomrule\end{longtable}
\begin{dosnotebox}\textbf{Bottom summary.} A gateway should help users depart toward their goal.\\[2mm]
\textbf{Your own synthesis:} \rule{.72\linewidth}{.4pt}\\[4mm]
\rule{.96\linewidth}{.4pt}\end{dosnotebox}

\section{Study Unit 09: D — Content Management}
\begin{longtable}{@{}p{.27\linewidth}p{.66\linewidth}@{}}\toprule
\textbf{Cue / Recall} & \textbf{Notes / Relationships}\\\midrule\endhead
How is reliable, reusable information produced? & Separate typed content from its rendering; establish editorial provenance, localization, update policy and search metadata.\\
\addlinespace[4mm]
\textbf{Practice} & Draft a content schema and a publishing-role permission matrix.\\\bottomrule\end{longtable}
\begin{dosnotebox}\textbf{Bottom summary.} Content operations are software and governance work.\\[2mm]
\textbf{Your own synthesis:} \rule{.72\linewidth}{.4pt}\\[4mm]
\rule{.96\linewidth}{.4pt}\end{dosnotebox}

\section{Study Unit 10: E — Trust and Credibility}
\begin{longtable}{@{}p{.27\linewidth}p{.66\linewidth}@{}}\toprule
\textbf{Cue / Recall} & \textbf{Notes / Relationships}\\\midrule\endhead
What evidence justifies customer trust? & Visual branding is not a substitute for transport security, accurate privacy notices, user control and phishing-resistant communication.\\
\addlinespace[4mm]
\textbf{Practice} & Walk through an account recovery threat model and consent revocation.\\\bottomrule\end{longtable}
\begin{dosnotebox}\textbf{Bottom summary.} Trust signals must correspond to real protections.\\[2mm]
\textbf{Your own synthesis:} \rule{.72\linewidth}{.4pt}\\[4mm]
\rule{.96\linewidth}{.4pt}\end{dosnotebox}

\section{Study Unit 11: F — Basic E-Commerce}
\begin{longtable}{@{}p{.27\linewidth}p{.66\linewidth}@{}}\toprule
\textbf{Cue / Recall} & \textbf{Notes / Relationships}\\\midrule\endhead
What can fail during purchase? & Ordering must handle concurrency, authoritative prices, secure payment boundaries, retries and explicit confirmation.\\
\addlinespace[4mm]
\textbf{Practice} & Design an idempotent checkout command and stock-change failure response.\\\bottomrule\end{longtable}
\begin{dosnotebox}\textbf{Bottom summary.} Transaction safety is part of UX.\\[2mm]
\textbf{Your own synthesis:} \rule{.72\linewidth}{.4pt}\\[4mm]
\rule{.96\linewidth}{.4pt}\end{dosnotebox}

\section{Study Unit 12: G — Advanced E-Commerce}
\begin{longtable}{@{}p{.27\linewidth}p{.66\linewidth}@{}}\toprule
\textbf{Cue / Recall} & \textbf{Notes / Relationships}\\\midrule\endhead
When does recommendation help versus harm? & Recommendation and cross-selling may improve discovery but introduce manipulation, fairness, review fraud and tracking risks.\\
\addlinespace[4mm]
\textbf{Practice} & Evaluate recommendations for cold start and opt-out, then review complaint handling.\\\bottomrule\end{longtable}
\begin{dosnotebox}\textbf{Bottom summary.} Relevance must not override agency.\\[2mm]
\textbf{Your own synthesis:} \rule{.72\linewidth}{.4pt}\\[4mm]
\rule{.96\linewidth}{.4pt}\end{dosnotebox}

\section{Study Unit 13: H — Task Completion}
\begin{longtable}{@{}p{.27\linewidth}p{.66\linewidth}@{}}\toprule
\textbf{Cue / Recall} & \textbf{Notes / Relationships}\\\midrule\endhead
How can people recover from mistakes and interruptions? & Multi-step flows need state machines, accessible feedback, safe persistence and authorized account operations.\\
\addlinespace[4mm]
\textbf{Practice} & Exercise a partially completed form through logout and network loss.\\\bottomrule\end{longtable}
\begin{dosnotebox}\textbf{Bottom summary.} Recovery is a first-class task requirement.\\[2mm]
\textbf{Your own synthesis:} \rule{.72\linewidth}{.4pt}\\[4mm]
\rule{.96\linewidth}{.4pt}\end{dosnotebox}

\section{Study Unit 14: I — Page Layout}
\begin{longtable}{@{}p{.27\linewidth}p{.66\linewidth}@{}}\toprule
\textbf{Cue / Recall} & \textbf{Notes / Relationships}\\\midrule\endhead
How do layout systems survive changing screens? & Semantic order, responsive reflow and legible information hierarchy matter more than a fixed desktop grid.\\
\addlinespace[4mm]
\textbf{Practice} & Review at 320px and 400 percent zoom with large localized strings.\\\bottomrule\end{longtable}
\begin{dosnotebox}\textbf{Bottom summary.} Structure should be resilient to user preferences.\\[2mm]
\textbf{Your own synthesis:} \rule{.72\linewidth}{.4pt}\\[4mm]
\rule{.96\linewidth}{.4pt}\end{dosnotebox}

\section{Study Unit 15: J — Site Search}
\begin{longtable}{@{}p{.27\linewidth}p{.66\linewidth}@{}}\toprule
\textbf{Cue / Recall} & \textbf{Notes / Relationships}\\\midrule\endhead
What makes results helpful and safe? & Search combines query entry, relevance ranking, index freshness, faceting and authorization filtering.\\
\addlinespace[4mm]
\textbf{Practice} & Build an evaluation corpus with zero-result and forbidden-document cases.\\\bottomrule\end{longtable}
\begin{dosnotebox}\textbf{Bottom summary.} Search is both navigation and data-access control.\\[2mm]
\textbf{Your own synthesis:} \rule{.72\linewidth}{.4pt}\\[4mm]
\rule{.96\linewidth}{.4pt}\end{dosnotebox}

\section{Study Unit 16: K — Navigation Usability}
\begin{longtable}{@{}p{.27\linewidth}p{.66\linewidth}@{}}\toprule
\textbf{Cue / Recall} & \textbf{Notes / Relationships}\\\midrule\endhead
How do controls communicate their purpose? & Links, buttons, tabs, breadcrumbs and error messages have distinct semantics; consistency and assistive-technology behavior are essential.\\
\addlinespace[4mm]
\textbf{Practice} & Compare a semantic link and an imperative click handler using keyboard access.\\\bottomrule\end{longtable}
\begin{dosnotebox}\textbf{Bottom summary.} Predictability comes from both words and behavior.\\[2mm]
\textbf{Your own synthesis:} \rule{.72\linewidth}{.4pt}\\[4mm]
\rule{.96\linewidth}{.4pt}\end{dosnotebox}

\section{Study Unit 17: L — Site Performance}
\begin{longtable}{@{}p{.27\linewidth}p{.66\linewidth}@{}}\toprule
\textbf{Cue / Recall} & \textbf{Notes / Relationships}\\\midrule\endhead
What should performance optimization measure? & Translate historical asset advice into user-centered loading, interaction and stability signals under contemporary transport and devices.\\
\addlinespace[4mm]
\textbf{Practice} & Instrument p75 LCP/INP/CLS; compare an image-priority optimization.\\\bottomrule\end{longtable}
\begin{dosnotebox}\textbf{Bottom summary.} Measure user impact rather than proxy counts alone.\\[2mm]
\textbf{Your own synthesis:} \rule{.72\linewidth}{.4pt}\\[4mm]
\rule{.96\linewidth}{.4pt}\end{dosnotebox}

\section{Study Unit 18: M — Mobile Web}
\begin{longtable}{@{}p{.27\linewidth}p{.66\linewidth}@{}}\toprule
\textbf{Cue / Recall} & \textbf{Notes / Relationships}\\\midrule\endhead
How do mobile needs alter the design? & Responsive layouts, touch and keyboard alternatives, intermittent connectivity and contextual device permissions redefine success.\\
\addlinespace[4mm]
\textbf{Practice} & Test denied geolocation and landscape/large-text behavior.\\\bottomrule\end{longtable}
\begin{dosnotebox}\textbf{Bottom summary.} Mobile access requires an inclusive fallback path.\\[2mm]
\textbf{Your own synthesis:} \rule{.72\linewidth}{.4pt}\\[4mm]
\rule{.96\linewidth}{.4pt}\end{dosnotebox}

\section{Active recall assessment}

\begin{enumerate}
\item Given a customer-facing feature, classify its source pattern IDs and the actual user outcome being optimized.
\item Defend why CSS Grid should replace table-layout practices while semantic data tables remain appropriate.
\item Explain why an accessible form can still fail a security requirement, and why a secure login can still be inaccessible.
\item Propose one observation, one negative test, one field metric and one rollback criterion for each critical user journey.
\item Explain when microfrontends add value, and why visual widget boundaries are not sufficient evidence for a service split.
\end{enumerate}


\clearpage
% design-of-sites-references.tex
% Shared references section for the Web Design Patterns / Design of Sites
% suite.
%
% RECONSTRUCTION PROVENANCE: the nine uploaded suite documents end with
% `\input{references.tex}`, but no `references.tex` was present in the upload
% batches or anywhere in this repository's history. The documents cite only
% keys [B1] and [B2] (bracketed tokens such as "[B4, F1, H10, ...]" are
% pattern-ID template placeholders, not citations). The two entries below are
% reconstructed from the bibliographic facts the documents themselves state:
% publisher record and bibliographic-database record of the second edition.
% This file is NOT the recovered original reference list.
\clearpage
\section*{References}
\begin{description}[leftmargin=3.2em,style=nextline]
  \item[{[B1]}] Publisher bibliographic record and public table of contents:
    D.~K.~van Duyne, J.~A.~Landay, J.~I.~Hong, \emph{The Design of Sites:
    Patterns, Principles, and Processes for Crafting a Customer-Centered Web
    Experience}, 2nd edition, Prentice Hall, December 2006, 1{,}024 pages.
    Confirms the five foundational chapters, the thirteen pattern groups
    (A--M), the 107 pattern labels and the research appendixes used as this
    suite's taxonomy.
  \item[{[B2]}] Bibliographic-database record (O'Reilly learning-platform
    catalog entry) for the same second edition; corroborates edition,
    authors, publication date and chapter/group structure. The copyrighted
    pattern explanations themselves were not retrieved or reproduced for
    this suite.
\end{description}

\end{document}
